\section{Introduction}
\label{sec:intro}

Self-supervised learning (SSL) on images has converged on a small set of
\emph{families}: contrastive instance discrimination
(SimCLR~\cite{chen2020simclr}), redundancy reduction
(VICReg~\cite{bardes2022vicreg}, Barlow Twins), predictive coding in latent
space (I-JEPA~\cite{assran2023ijepa}), and recent Gaussian-regularized JEPA
variants (LeJEPA~\cite{balestriero2025lejepa}).
Each family optimizes a different geometry: InfoNCE excels at local
neighborhood structure (k-NN), while spread and invariance terms often
improve linear separability.

We ask a practical question for low-resource benchmarks: \emph{given a fixed
$\sim$3M-parameter encoder and a masked JEPA view pipeline, how do published
objectives compare at long training (200 epochs), and can we combine the
strengths of contrastive mask prediction with non-contrastive spread?}

Our contributions are:
\begin{enumerate}
  \item A unified CIFAR-100 benchmark with matched architecture, corruption,
        augmentation, and evaluation (Section~\ref{sec:evaluation}).
  \item An ablation over representative published losses---JEPA cosine,
        InfoNCE, VICReg, SIGReg, uniformity, and SwAV---at 200 epochs on the
        same SmallResNet (Section~\ref{sec:ablation}).
  \item \textbf{Asymmetric hybrids} that keep InfoNCE on the mask/predictor
        path and replace augmentation InfoNCE with MSE+var+cov or MSE+SIGReg
        (Section~\ref{sec:method}), improving linear probe over dual InfoNCE
        with a controllable k-NN cost.
\end{enumerate}
